% ------------------------------------------------------------
%  Preamble for the standalone paper.  Deliberately independent
%  of annotated.sty: standard article apparatus, normal margins,
%  no margin-note column, no house boxes.  Only the TikZ drawing
%  vocabulary is carried over, so the figures survive unchanged.
% ------------------------------------------------------------
\usepackage[T1]{fontenc}
\usepackage{amsmath,amssymb,amsthm,mathtools}
\usepackage{mathrsfs}
\usepackage[margin=1.15in]{geometry}
\usepackage{xcolor}
\usepackage{tikz}
\usepackage{booktabs}
\usepackage{enumitem}
\usepackage{microtype}
\usepackage[hidelinks,hypertexnames=false]{hyperref}
\usetikzlibrary{arrows.meta,positioning,patterns}

\definecolor{RuInk} {gray}{0.00}
\definecolor{RuGrey}{gray}{0.42}

\tikzset{
  axis/.style   ={draw=RuInk, line width=0.5pt, -{Stealth[length=4pt]}},
  seg/.style    ={draw=RuInk, line width=1.2pt},
  thin/.style   ={draw=RuInk, line width=0.4pt},
  guide/.style  ={draw=RuInk, line width=0.3pt,
                  dash pattern=on 1.2pt off 1.2pt},
  ptfill/.style ={circle, fill=RuInk, inner sep=0pt, minimum size=3.2pt},
  ptopen/.style ={circle, draw=RuInk, line width=0.5pt, fill=white,
                  inner sep=0pt, minimum size=3.4pt},
  lbl/.style    ={font=\scriptsize\sffamily, inner sep=1.5pt},
  mlbl/.style   ={font=\scriptsize, inner sep=1.5pt},
  rmbox/.style  ={draw=RuInk, line width=0.5pt, fill=white,
                  minimum height=11mm, inner sep=4pt, align=center,
                  font=\footnotesize\sffamily},
  rmarrow/.style={-{Stealth[length=4.5pt]}, draw=RuInk, line width=0.5pt}}

\theoremstyle{plain}
\newtheorem{theorem}{Theorem}[section]
\newtheorem{proposition}[theorem]{Proposition}
\newtheorem{lemma}[theorem]{Lemma}
\newtheorem{corollary}[theorem]{Corollary}
\theoremstyle{definition}
\newtheorem{definition}[theorem]{Definition}
\newtheorem{example}[theorem]{Example}
\theoremstyle{remark}
\newtheorem{remark}[theorem]{Remark}

% A short strategy paragraph before a substantial proof.  Kept as plain
% typography rather than a boxed environment: this is a paper, not the
% annotated edition.
\newcommand{\idea}[1]{\par\smallskip\noindent
  \textit{Idea of the proof.}\enspace #1\par\smallskip}

% A numbered walk through the architecture of a proof, in plain
% typography.  Used for the pillar proofs; \idea for the shorter ones.
\newenvironment{flow}[1][How the proof runs]{%
  \par\smallskip\noindent\textit{#1}%
  \begin{enumerate}[leftmargin=2.1em, itemsep=2.5pt, topsep=3pt,
                    label=\textup{(\arabic*)}]}
  {\end{enumerate}\par\smallskip}

\newcommand{\R}{\mathbb{R}}
\newcommand{\Q}{\mathbb{Q}}
\newcommand{\N}{\mathbb{N}}
\newcommand{\1}{\mathbf{1}}
\emergencystretch=1.2em
